\documentclass[
    language=english,
    doctype=legal-brief,
    institution=none,
]{../../omnilatex}

\RequirePackage{config/document-settings}

\begin{document}
\maketitle

\section*{IN THE UNITED STATES DISTRICT COURT FOR THE NORTHERN DISTRICT OF CALIFORNIA}

\noindent\textbf{Case No.:} 3:24-cv-01847-WHO\\
\textbf{Judge:} William H. Orrick III

\vspace{0.5cm}

\noindent\textbf{INNOVATE TECH CORP.,} a Delaware corporation,\\
\hspace{1.5cm}\textit{Plaintiff,}\\[0.3cm]
\textbf{v.}\\[0.3cm]
\textbf{APEX DIGITAL SOLUTIONS, INC.,} a Nevada corporation,\\
\hspace{1.5cm}\textit{Defendant.}\\[0.5cm]

\noindent\rule{\textwidth}{0.4pt}\\[0.3cm]

\noindent\textbf{PLAINTIFF'S MEMORANDUM OF POINTS AND AUTHORITIES IN OPPOSITION TO DEFENDANT'S MOTION TO DISMISS}

\vspace{0.5cm}

\section{QUESTIONS PRESENTED}

\begin{enumerate}
    \item Whether Plaintiff has stated a plausible claim for trade secret misappropriation under the Defend Trade Secrets Act, 18 U.S.C. \S 1836, where Defendant accessed Plaintiff's proprietary machine learning algorithms through a former employee who had signed a non-disclosure agreement and subsequently joined Defendant's competing division.

    \item Whether Plaintiff's claim for breach of the duty of loyalty survives a motion to dismiss where the former employee solicited Plaintiff's key technical staff and directed them to download proprietary source code to personal devices before departing to Defendant.

    \item Whether Defendant's conduct constitutes unfair competition under California Business and Professions Code \S 17200, given the pattern of targeted hiring combined with systematic extraction of trade secrets.
\end{enumerate}

\section{STATEMENT OF THE CASE}

Plaintiff Innovate Tech Corp. (``Innovate'' or ``Plaintiff'') brings this action against Defendant Apex Digital Solutions, Inc. (``Apex'' or ``Defendant'') for the willful and deliberate misappropriation of trade secrets, breach of contractual and fiduciary duties, and unfair business practices arising from Defendant's systematic campaign to acquire Plaintiff's proprietary technology through the recruitment of Plaintiff's senior engineers.

\subsection{Plaintiff's Proprietary Technology}
Plaintiff has invested over \$45 million and eight years of research and development in creating its proprietary machine learning platform, ``NeuralForge.'' NeuralForge represents a significant competitive advantage in the enterprise analytics market, enabling real-time predictive modeling with 40\% greater accuracy than industry benchmarks. The platform's core algorithms, training data pipelines, and architectural innovations constitute valuable trade secrets.

\subsection{The Departure and Misappropriation}
Between January and March 2024, four senior engineers from Plaintiff's NeuralForge team resigned and accepted positions at Apex's newly created ``Advanced Analytics Division.'' Each engineer had executed Confidentiality and Invention Assignment Agreements that explicitly prohibited the use or disclosure of Plaintiff's trade secrets.

Evidence obtained through forensic analysis of company devices reveals that, prior to their departure:
\begin{enumerate}[label=(\alph*)]
    \item The engineers collectively downloaded over 12,000 files containing proprietary source code, model weights, and training data specifications;
    \item These downloads were made to personal devices and external storage media;
    \item The engineers accessed documents and repositories for which they had no legitimate business need in the weeks preceding their departure;
    \item Communications recovered from company email show coordination with Apex personnel regarding the acquisition of Plaintiff's technology.
\end{enumerate}

\subsection{Procedural History}
Plaintiff filed its Complaint on April 15, 2024, asserting causes of action for: (1) trade secret misappropriation under the DTSA; (2) trade secret misappropriation under the California Uniform Trade Secrets Act; (3) breach of contract; (4) breach of the duty of loyalty; and (5) unfair competition under \S 17200. Defendant filed the instant Motion to Dismiss on May 30, 2024, pursuant to Federal Rule of Civil Procedure 12(b)(6).

\section{LEGAL STANDARD}

On a motion to dismiss under Rule 12(b)(6), the court must accept the well-pleaded factual allegations in the complaint as true and draw all reasonable inferences in the plaintiff's favor. \textit{Ashcroft v. Iqbal}, 556 U.S. 662, 678 (2009). A claim is facially plausible when the factual content allows the court to draw a reasonable inference that the defendant is liable for the misconduct alleged. \textit{Id.} at 678.

For trade secret claims, the complaint must identify the trade secrets with reasonable specificity and allege facts supporting the reasonable measures taken to maintain their secrecy. \textit{MagID Security, Inc. v. Cardio Security International, Inc.}, 224 F. Supp. 2d 1222, 1227 (N.D. Cal. 2002).

\section{ARGUMENT}

\subsection{Plaintiff Adequately Alleges Trade Secret Misappropriation}

Defendant contends that Plaintiff fails to identify its trade secrets with sufficient specificity. This argument is without merit. The Complaint identifies the trade secrets with particularity, including:

\begin{enumerate}
    \item The NeuralForge predictive modeling algorithms, including the proprietary attention mechanism architecture described in Exhibits B-1 through B-4;
    \item The training data curation and preprocessing pipelines, including the proprietary data augmentation techniques;
    \item The model optimization procedures that achieve superior performance metrics;
    \item The system architecture and deployment configurations.
\end{enumerate}

Courts in this District have consistently held that identifying trade secrets by reference to specific technical documents, code repositories, and architectural descriptions satisfies the specificity requirement. \textit{See Palantir Technologies Inc. v. GoodData Corp.}, No. 15-cv-03732, 2016 WL 3345688 (N.D. Cal. June 15, 2016) (upholding trade secret identification where plaintiff referenced specific source code modules and algorithms).

Moreover, Plaintiff has demonstrated that it took reasonable measures to protect these secrets, including: executing confidentiality agreements with all employees with access to the technology; implementing access controls limiting access to need-to-know personnel; using encrypted repositories and version control systems; conducting exit interviews and reminding departing employees of their obligations; and monitoring for unauthorized downloads and data exfiltration.

\subsection{Plaintiff's Breach of Duty of Loyalty Claim Is Properly Pled}

Defendant argues that the duty of loyalty claim fails because the engineers were not agents of Defendant at the time of the alleged misconduct. Defendant misconstrues the law. The duty of loyalty applies to the engineers themselves as employees of Plaintiff, and Defendant's liability arises from its knowing participation in and inducement of the breach.

Under California law, an employer may bring a claim for breach of the duty of loyalty against a former employee who, during the employment relationship, engaged in conduct adverse to the employer's interests, including soliciting employees to leave and misappropriating trade secrets. \textit{See Eastman Kodak Co. v. Banner Office Supplies, Inc.}, 238 Cal. App. 3d 1591, 1602 (1991).

Furthermore, Defendant's liability is independently established through the doctrine of tortious interference. Defendant knowingly induced the engineers to breach their confidentiality agreements and duty of loyalty. \textit{See Reagan v. Regents of the University of California}, 194 Cal. App. 3d 555, 564 (1987).

\subsection{Plaintiff States a Plausible \S 17200 Claim}

Defendant argues that Plaintiff's \S 17200 claim is duplicative of its trade secret claims. This argument fails because \S 17200 operates as an independent statutory cause of action that does not require allegations of actual competition or an actionable underlying wrong.

The California Supreme Court has held that ``an allegation of an underlying wrongdoing is not a prerequisite to stating a claim'' under \S 17200. \textit{Kearns v. Ford Motor Co.}, 56 Cal. 4th 132, 142 (2013). Here, Plaintiff's allegations that Defendant engaged in a deliberate campaign to acquire trade secrets through targeted recruitment and systematic data exfiltration states a claim for ``any unlawful, unfair, or fraudulent business act or practice.'' Cal. Bus. \& Prof. Code \S 17200.

Defendant's conduct constitutes ``unfair'' competition because: (1) the injury to consumers and competitors outweighs any business justification; (2) the conduct is immoral, unethical, oppressive, or harmful; and (3) it threatens fair competition. \textit{See Cel-Tech Communications, Inc. v. Los Angeles Cellular Telephone Co.}, 20 Cal. 4th 163, 186 (1999).

\section{CONCLUSION}

For the foregoing reasons, Plaintiff respectfully requests that Defendant's Motion to Dismiss be denied in its entirety.

\vspace{2cm}

\noindent\rule{6cm}{0.4pt}\\
\textbf{JENNIFER A. MARSHALL, ESQ.}\\
Cal. Bar No. 245891\\
Marshall \& Associates LLP\\
100 Montgomery Street, Suite 2400\\
San Francisco, CA 94104\\
Telephone: (415) 555-0188\\
Email: \url{jmarshall@marshallassociates.com}

\vspace{1cm}

\noindent\textit{Attorneys for Plaintiff}\\
\textbf{INNOVATE TECH CORP.}

\end{document}
